\section{Introduction}
\label{sec:introduction}

\subsection{Purpose of the Document}

This document defines the proposed requirements
and architecture of OliveTrack, a platform
for managing olive harvest operations in a private deployment for Pepe's parents,
with one initial owner account and owner-invited managers.

It establishes a shared understanding of the problem,
intended users, functional scope, business rules
and expected system behaviour.
It also connects the use cases to the functional
requirements, quality requirements
and architectural decisions developed by the team.
Section~\ref{sec:risk-management} also records the working delivery-risk assessment
and response approach. The concise summary in Section~\ref{sec:project-planning}
connects the specification to effort, delivery planning and progress evidence.

The document supports both product development
and project management.
An explicit description of the expected behaviour
helps the team identify missing work,
define implementation and validation tasks,
allocate responsibilities and prepare effort estimates.

This working specification complements the estimation workbook, team hours
sheet and other course deliverables. It provides a shared requirements and design
baseline for review; it does not replace task assignments, an approved delivery
plan or acceptance evidence.

\subsection{Project Context}

Olive harvest management involves information
that is produced by different people
and becomes available at different times.

Before fieldwork, the owner may optionally plan the expected workers,
work areas and, where known, schedules and tasks.
During or after the working day, actual attendance,
work periods, tasks and breaks must be recorded.

Harvested olives are transported and delivered
to a cooperative or another receiving destination.
Weighing tickets provide information about
the delivered load.
Analysis results and monetary values may arrive later.

Worker documents, including contracts
and insurance records, must also be accessible.
Some documents require a signing interaction.
Insurance information is particularly important
when work takes place in the field.

When this information is distributed across
paper records and separate files,
it becomes harder to retrieve, reconcile
and summarise the campaign's results.
Differences between planned and actual work,
missing analysis results and ambiguous weight values
can also lead to incorrect conclusions.

OliveTrack proposes a shared digital environment
for the owner's business.
It supports properties with optional parcels,
harvest campaigns covering one or more areas,
worker records, operational entry,
document consultation and campaign indicators.

Field connectivity may be unreliable.
Essential field operations therefore require
prepared offline access and local entry,
followed by synchronisation when connectivity returns.
Administrative operations may require connectivity,
as specified in the individual use cases.

The project is being developed within the
Project Management course.
The proposed scope must remain compatible
with the academic schedule,
available resources and agreed delivery objectives.

\subsection{Objectives}

The proposed solution pursues the following objectives:

\begin{enumerate}

    \item \textbf{Organise the harvest context.}
    Maintain identifiable properties, optional parcels
    and campaigns without mixing operational records
    from different areas or harvest periods.

    \item \textbf{Separate planning from actual work.}
    Support daily and multi-day planning,
    while requiring explicit confirmation
    of attendance and actual work.

    \item \textbf{Record work consistently.}
    Distinguish machinery operation from general work,
    account for explicitly recorded breaks
    and prevent overlapping actual periods
    for the same worker.

    \item \textbf{Track deliveries and source information.}
    Connect transport, delivery, weighing,
    analysis and monetary information
    without inventing missing measurements
    or applying deductions twice.

    \item \textbf{Support worker-document workflows.}
    Provide controlled access to uploaded documents,
    prepared offline insurance consultation
    and the agreed signing interaction.

    \item \textbf{Provide understandable indicators.}
    Present campaign results with clear units,
    reporting scope, data completeness
    and provisional or final status where applicable.

    \item \textbf{Support work in the field.}
    Allow essential operations to continue
    with prepared authorised data
    when a network connection is unavailable.


\end{enumerate}

Detailed measurable targets are addressed
in the non-functional requirements.
The objectives above do not constitute
verified performance or operational results.

\subsection{Sources and Document Status}\label{sec:document-status}

The current document draws on:

\begin{itemize}

    \item The \emph{OliveTrack Team Guide},
    dated 9 October 2026~\cite{project-guide},
    describing the business context,
    illustrative examples and unresolved questions.

    \item The existing Statement of Work,
    persona descriptions, general architecture,
    data architecture and non-functional
    requirements proposals supplied to the team~\cite{project-baseline}.

    \item Initial responses supplied
    by the teammate acting as an intermediary
    with the client.

    \item The subsequent requirements review,
    including the use-case discussions
    and proposed resolutions recorded
    during document development.

    \item The team's security, privacy and reliability review and the updated
    performance, usability and availability proposal dated 3 November 2026,
    reconciled with the current functional decisions.

    \item The team's \emph{OliveTrack General Architecture}, version 1.1, and
    \emph{OliveTrack Risk Management Plan}, dated November 2026~\cite{project-architecture,project-risk}, reconciled with
    the current functional/data baseline. Unconfirmed voting, milestone, resource
    and provider assumptions retain explicit review status.

    \item Operational photographs and ticket explanations supplied by the client
    intermediary on 6 November 2026, reviewed on 21 November 2026~\cite{project-examples}. Field interpretation
    and form inventory appear in Section~\ref{sec:source-examples}.
\end{itemize}

This working specification consolidates the supplied sources and subsequent
requirements, business-rule and technical-decision reviews, including the
stack, access model and quality targets adopted during the author's review of
5 November 2026, followed by the operational-document review of 21 November 2026. Illustrative examples, working decisions and client-provided
responses have different validation status. Conflicting
statements require an explicit decision; author review does not establish formal
approval by the whole team or client, or legal validation of templates or signatures.

Section~\ref{sec:validation} records current decisions, residual questions and
validation responsibilities; the detailed questionnaire is maintained separately.
The intermediate submission milestone was 6 November 2026; the presentation
milestone was 10 November 2026, with a browser-based navigable prototype and sample
data as the planned demonstration. This revision is dated 21 November 2026;
the milestone dates alone do not establish submission or presentation evidence.
Throughout this report, specified behaviour, prototype interaction, implemented
functionality and verified results are distinct. The demonstration does not
establish production security, durable synchronisation or operational acceptance.

\paragraph{Project Deliverables.}
The files developed during this project phase are collected in the shared project
folder: this report and its LaTeX source, effort estimates, recorded team hours,
the supplementary data reference, team roster and navigable demonstration.
The presentation slides were planned for the 10 November 2026 milestone;
their delivery status must be checked in the shared folder.
\par\href{https://drive.google.com/drive/folders/1rIQ5pGJDSKda0syKZ2MtbdDA0DTqOnRW?usp=sharing}{Open the shared project folder.}

\paragraph{How to use this specification.}
Sections~\ref{sec:users} and~\ref{sec:scope} establish stakeholders and boundaries;
Sections~\ref{sec:use-cases}--\ref{sec:non-functional-requirements} define behaviour,
shared constraints and quality criteria. Sections~\ref{sec:solution-architecture}
and~\ref{sec:data-architecture} describe the solution and data design.
Section~\ref{sec:risk-management} connects risks to delivery planning, and
Section~\ref{sec:validation} records validation, decisions and unresolved matters. Read the full data dictionary only when
implementing or checking fields not repeated in the principal model.

\paragraph{Course alignment and AI assistance.}
The supplied course materials guide requirement organisation/change control,
risk assessment and planning~\cite{course-presentation,course-concepts,course-scope,course-estimation,course-risk,course-planning,course-customer}. Templates are used
selectively, without claiming that every taught process has been performed.
AI assistance was used in drafting/restructuring this report, preparing LaTeX
and developing the navigable prototype. These contributions require author
verification; they are not client approval, executed acceptance tests or evidence
of production readiness. The team-wide inventory of tools, contributions and
review responsibilities remains to be completed under \decisionref{Q-PM-02}
in accordance with the course disclosure requirements.

\paragraph{Revision of 24 November 2026.}
The owner confirmed private initial provisioning by a Technical Administrator;
manager account creation remains invitation-based and planning remains included.
The attached skeleton, hours sheet, estimation workbook and revised task catalogue
supply implementation/planning evidence, with their different baselines identified
in Section~\ref{sec:project-planning}. Existing scaffold files are evidence of a
starting structure; executed tests are needed for completion claims.
